\section{A spatial measure with dimension greater than seven quarters}
\label{sec:spatial}

The boundary argument will count participating vertices, so it needs a
strip estimate for vertex mass while retaining a positive amount of edge
mass. We now construct a single Euclidean limit with these two properties.
The metric estimate from the preceding section concerns one fixed rank
point and radius at a time. A Hausdorff limit of graphs will allow us to
use it with configuration thresholds that still depend on the point
and radius.

Write $D=(0,1)^2$, and let $\pi:Z\to[0,1]^2$ be the projection onto the
rank coordinates $(z_A,z_B)$.  Denote the inverse rank maps, in the common
Euclidean normalization, by $L_n:D\to\mathbb R^2$.  These are
homeomorphisms.  For $\overline B(p,h)\subset D$, put
\[
 U_n(p,h)=\max_{|u-p|=h}|L_n(u)-L_n(p)|,
 \qquad
 V_n(p,h)=\min_{|u-p|=h}|L_n(u)-L_n(p)|.
\]
The input from the preceding section has the following exact quantifiers.
For each selected $p$, there is $C_p<\infty$ such that, for every integer
$T\ge2$, all sufficiently small fixed $h>0$ satisfy, eventually in $n$,
\begin{equation}
 V_n(p,h/T)\ge\frac{U_n(p,h)}{C_pT},
 \qquad
 U_n(p,h/T)\le\frac{C_pV_n(p,h)}{T}.
 \label{eq:spatial-metric-input}
\end{equation}
The constant $C_p$ is independent of $T$.  The allowed radii may depend on
$p,T$, and the eventual index may depend on $p,T,h$.

We use the positive good part already selected above: a Borel set
$E\subset Z$ with $\lambda_{K_0}(E)>0$, whose restricted endpoint measure
has absolutely continuous pair-coordinate projection, and for which
\eqref{eq:spatial-metric-input} holds at every $p\in\pi(E)$ after deleting
a null set.  We also retain $\lambda\ll\mu$.

The absolutely continuous projection here belongs to edge mass
$\lambda|_E$, not yet to the vertex mass $\mu|_E$ whose strip mass we
must bound. We will discard labels over a null set carrying the singular
part of the vertex projection and restrict its density to a bounded
range. The first deletion loses no selected edge mass, by absolute
continuity of its projection, and some finite density bound still
retains positive $\lambda_{K_0}$-mass. The metric construction will
then transfer this vertex-density bound from
rank coordinates to Euclidean disks and strips.

\begin{proposition}\label{prop:spatial-limit}
There is a Borel set $G\subset E$ with $\lambda_{K_0}(G)>0$, a further
subsequence, and a translation and uniform rescaling of each normalized
Euclidean plane with the following properties.  Over every
$p\in\pi(G)$ the limiting graph of the maps $L_n$ has a unique position
$l(p)$ in the closed unit disk.  With $\eta=1/16$, there is $c_G>0$ such
that
\begin{equation}
 |l(p)-l(q)|\ge c_G|p-q|^{1+\eta}
 \qquad(p,q\in\pi(G)).
 \label{eq:inverse-holder}
\end{equation}
Attach Euclidean position as an additional compactified coordinate to the
vertex measures, and take a weak limit $\widetilde\mu$.  Its restriction
by the old-label condition $z\in G$ has position marginal $\nu_G$
satisfying
\begin{equation}
 \nu_G(B(x,r))\le C_G r^{d_*},
 \qquad d_*=\frac{2}{1+\eta}=\frac{32}{17},
 \qquad x\in\mathbb R^2,\quad 0<r<1.
 \label{eq:disk-bound}
\end{equation}
Consequently every strip of width $r$ has $\nu_G$-mass at most
$C'_G r^{d_*-1}$, uniformly in its position and direction.  Every line
has zero $\nu_G$-mass.
\end{proposition}

\begin{proof}
We first select a positive set with common scale parameters and bounded
vertex density in the rank coordinates. After normalizing each
configuration at one common rank point, upper radial bounds will make the limiting graph single-valued
on that set, and lower radial bounds will give \eqref{eq:inverse-holder}.
Combining this estimate with the vertex density will prove the disk
and strip bounds.

For each integer $T\ge2$ and integer $m\ge0$, set
$h_j=T^{-m-j}$ and consider the points for which
$\overline B(p,h_0)\subset D$ and, for every $j\ge0$, eventually in $n$,
\begin{equation}
 \begin{aligned}
 V_n(p,h_{j+1})&\ge \alpha U_n(p,h_j),
 &\alpha&=T^{-1-\eta},\\
 U_n(p,h_{j+1})&\le \beta V_n(p,h_j),
 &\beta&=T^{-1+\eta}.
 \end{aligned}
 \label{eq:scale-radii}
\end{equation}
These are Borel conditions.  Indeed, at a fixed admissible radius,
$U_n$ and $V_n$ are continuous functions of the center, and eventual
validity is a countable union over starting indices followed by a
countable intersection over later indices.  Taking the further countable
intersection over $j$ preserves measurability.

For each point satisfying \eqref{eq:spatial-metric-input}, choose $T$
with $C_p\le T^\eta$, and then choose $m$ large enough.  This point
satisfies \eqref{eq:scale-radii} at all the selected radii.  The countable
union over $T,m$ therefore covers the selected good part.  Fix one pair
$T,m$ retaining positive $\lambda_{K_0}$-mass.  This selection makes no
claim of a common starting index.

The bounded density must concern vertices, since their activity will be
estimated later.  Decompose the pair-coordinate projection of the entire
vertex measure as
\[
 \pi_\#\mu=f(p)\,dp+\sigma_s,
 \qquad \sigma_s\perp dp.
\]
Choose an area-zero Borel carrier $S$ for $\sigma_s$, enlarged by any
area-zero set on which $f$ is not finite.  Delete $\pi^{-1}(S)$ from the
selected labels.  This loses no good mass because their endpoint
projection is absolutely continuous.  Next restrict to
$f(\pi(z))\le M$ for a sufficiently large finite $M$, retaining positive
$\lambda_{K_0}$-mass. Denote the resulting set by $G$; it satisfies
\begin{equation}
 \mu\bigl(G\cap\pi^{-1}(B)\bigr)\le M|B|
 \qquad\text{for every Borel }B\subset D.
 \label{eq:vertex-density}
\end{equation}
Here we have removed the singular part of the vertex projection and
truncated its density; absolute continuity of the endpoint projection
alone would not imply \eqref{eq:vertex-density}.

Cover $D$ by countably many balls of radius $h_0/8$ and restrict $G$ to
one of positive $\lambda_{K_0}$-mass.  Choose $p_0$ in its projected good
set.  Then
\[
 |p-p_0|<h_0/4\qquad(p\in\pi(G)).
\]
Translate and uniformly rescale the maps so that
\[
 L_n(p_0)=0,\qquad U_n(p_0,h_0)=1.
\]
The scale factor is positive, and \eqref{eq:scale-radii} is unchanged.

We next obtain bounds at every fixed scale, with constants independent of
the center.  For a homeomorphism, the image of a closed disk is a
topological closed disk whose boundary is the image of the boundary
circle.  Thus
\begin{equation}
 B(L_n(p),V_n(p,h))\subset L_n(B(p,h))
 \subset \overline B(L_n(p),U_n(p,h)).
 \label{eq:radial-enclosure}
\end{equation}
For the first inclusion, a segment from $L_n(p)$ to a missing point
closer than $V_n(p,h)$ would first leave the image through a closer
boundary point.  For the second, a positive maximum of the distance from
$L_n(p)$ cannot occur in the interior of the image disk.  Also, a point
outside $B(p,h)$ has image distance at least $V_n(p,h)$ from $L_n(p)$.

Since $T\ge2$, the localization gives
\[
 B(p_0,h_1)\subset B(p,h_0),
 \qquad
 \overline B(p,h_1)\subset B(p_0,h_0).
\]
Eventually $V_n(p_0,h_1)\ge\alpha$, so the first image contains $B(0,\alpha)$.
Its maximum distance from the possibly nonzero center $L_n(p)$ is at
least $\alpha$.  The second image is contained in the unit disk.  Therefore
\[
 U_n(p,h_0)\ge\alpha,\qquad U_n(p,h_1)\le2,
 \qquad |L_n(p)|\le1.
\]
Only the first inequality needs the anchor's eventual index.  The factor
$2$ occurs because $U_n(p,h_1)$ is measured from $L_n(p)$, rather than
from the origin.

Fix $p\in\pi(G)$ and an integer $j\ge1$.  Iterating the finitely many
inequalities \eqref{eq:scale-radii}, and using $V_n\le U_n$, gives,
eventually in $n$,
\begin{equation}
 V_n(p,h_j)\ge\alpha^{j+1},
 \qquad U_n(p,h_j)\le2\beta^{j-1}.
 \label{eq:iterated-radii}
\end{equation}
The constants are uniform over $G$.  The index threshold may depend on
$p,j$.  We will always fix these parameters before taking the limit.

The upper bounds in \eqref{eq:iterated-radii} tend to zero with the
scale, while the lower bounds keep distinct rank points separated.
We now pass to one graph limit in which both facts can be used, keeping
each center and scale fixed before taking the configuration limit.

To put all spatial limits on one subsequence, let $X$ be a compact metric
compactification of $\mathbb R^2$, and take the closures of the graphs
of $L_n$ in $[0,1]^2\times X$.  Pass to a subsequence on which these
compact sets converge in Hausdorff distance to a closed relation
$\Gamma$.  At an interior rank point, each finite-index closed graph
contains only the actual graph point, by continuity of $L_n$.

Fix $p\in\pi(G)$.  Its limiting fiber is nonempty by compactness,
because $|L_n(p)|\le1$.  Every $(p,x)\in\Gamma$ has graph-point
approximations $(u_n,L_n(u_n))$ at every sufficiently late index of the
same subsequence.  Since $u_n\to p$, eventually
$u_n\in B(p_0,h_0)$; consequently $|L_n(u_n)|\le1$.
Thus all points in this fiber have finite Euclidean positions.

Suppose $x$ and $y$ are two such positions.  Hausdorff convergence
provides approximations for both at the same indices.  Fix $j\ge1$.
Both rank arguments eventually lie in $B(p,h_j)$, where
\eqref{eq:radial-enclosure} and \eqref{eq:iterated-radii} bound the
distance between their positions by $4\beta^{j-1}$.  Passing to the limit
and then letting $j\to\infty$ proves $x=y$, since $\beta<1$.
Call this point $l(p)$.  Uniqueness also shows that $L_n(p)\to l(p)$
along the full selected subsequence.

For distinct $p,q\in\pi(G)$ put $d=|p-q|<h_0/2$.  Choose the first
$j\ge1$ with $h_j<d$; then $h_j\ge d/T$.  The point $q$ lies outside
the radius-$h_j$ disk centered at $p$, so eventually
\[
 |L_n(q)-L_n(p)|\ge V_n(p,h_j)\ge\alpha^{j+1}.
\]
Both fixed-point sequences converge on our subsequence.  Therefore
\[
 |l(q)-l(p)|\ge
 \alpha\left(\frac{h_j}{h_0}\right)^{1+\eta}
 \ge \frac{\alpha}{(Th_0)^{1+\eta}}\,d^{1+\eta},
\]
which proves \eqref{eq:inverse-holder}.  This argument uses a finite
number of scales for each fixed pair, not a common threshold over $G$.

Finally attach compactified position to the vertex and other spatial
measures before taking weak limits on this same subsequence.  Their
old-coordinate projections are unchanged.  Their pair-position
projections are supported on $\Gamma$: any compact set disjoint from
$\Gamma$ is disjoint from every sufficiently late closed graph, so
continuous tests supported there have limiting integral zero.
Consequently the enhanced vertex measure restricted by $z\in G$ lies
over the unique positions $l(\pi(z))$.  This restriction is made after
taking the limit; no convergence of finite-index restrictions to an
arbitrary Borel set $G$ is asserted.

If a disk of radius $r$ contains positions over $p$ and $q$, then
\eqref{eq:inverse-holder} gives
\[
 |p-q|\le (2r/c_G)^{1/(1+\eta)}.
\]
Thus its pair-coordinate preimage has planar outer area at most a
constant times $r^{2/(1+\eta)}$.  Applying
\eqref{eq:vertex-density} to that preimage proves
\eqref{eq:disk-bound}.  Using outer area avoids any requirement that a
projection of a Borel set itself be Borel.

The support of $\nu_G$ lies in the unit disk.  Its intersection with a
strip of width $r$ is covered by $O(1/r)$ disks of radius $O(r)$.
Summing \eqref{eq:disk-bound} proves the strip estimate; shrinking strips
to a line proves that lines have zero mass.
\end{proof}

The same enhancement preserves $\widetilde\mu_J\le\widetilde\mu$ and
all old marginals, in particular the positive mass
$\widetilde\lambda_{K_0}\{z\in G\}=\lambda_{K_0}(G)$.
By the enhanced-coordinate form of the earlier crossing inequality,
that inequality is also available on this selected set.  We will use the strip estimate
for its active-vertex measures.  The numerical fact needed below is
\[
 d_*=\frac{32}{17}>\frac74.
\]
This is an auxiliary spatial exponent; it is not the exponent saving in
the final halving-line bound.
